% sections/theory.tex --- Minos manuscript, theory section.
%
% Standalone section for later \input into the Minos manuscript. Compiles via the harness
% paper/theory_standalone.tex (CP2). The including document must provide, in its preamble:
%   \usepackage{amsmath,amssymb,booktabs,microtype}
%   \usepackage{amsthm}                              % theorem/proposition/proof environments
%   \newtheorem{theorem}{Theorem}  \newtheorem{proposition}{Proposition}
%   \usepackage[numbers,sort&compress]{natbib}       % numbered (Vancouver) citations
%   % Overleaf font (standing convention): \usepackage{ebgaramond}  -- commented for local compile.
% Notation matches the locked model in theory/THEORY_MODEL.md.
%
% PROVENANCE. Every numeric value below carries a "% src:" comment naming the script and seed that
% produced it, all re-run this session in the `proteus` env. Theory constants are sympy-derived
% (theory/gap_scaling.py); the bound constants from theory/detectability.py; the theory-vs-simulation
% confirmation from theory/confirm.py; v2/v3 simulation references from minos-core/RESULTS_B.md
% (seed 20240517, RUN_N = 4e6) and minos-core/RESULTS_C.md (CAL_N = 2e6, DEP_N = 1e6). The
% impossibility (Theorem~\ref{thm:detect}(i)) is a completed data-processing proof, machine-verified
% in theory/impossibility_check.py (GATE 2(i)); it is conditional only on the definition of "hidden"
% as the observable-invariant component of a shift.

\section{Theory: the decision--calibration gap and the limits of label-free validity monitoring}
\label{sec:theory}

A reported error bar can be optimal for two different jobs, and the two need not agree. We make this
precise for the Minos model of a piecewise-linear clinical utility acting on a skewed posterior, and
prove two statements about it. Theorem~\ref{thm:gap} turns the empirically observed
\emph{decision--calibration gap} into a leading-order scaling law in the posterior skew, with a
coefficient that is a pure decision-boundary quantity. Theorem~\ref{thm:detect} splits the problem of
monitoring a deployed correction for staleness into an observable component that is
bounded-detectable and a hidden component that is undetectable by \emph{any} label-free statistic.
The latter is the formal case for sparse labeled spot-checks, and it is an information-theoretic
impossibility distinct in kind from an identifiability limit.

\subsection{Setup: two calibrations of one error bar}
\label{sec:setup}

Following the locked model (the analytic transcription of the simulation code; see
\texttt{theory/\allowbreak THEORY\_MODEL.md}), reports are generated posterior-centrically: the reported centre is
$\mu$, and the truth is $\theta = \mu + s\,u$, where $u$ is a zero-mean, unit-variance skew-normal
error of shape $\kappa$. Its standardised third cumulant is
\begin{equation}
\gamma(\kappa) \;=\; \tfrac{4-\pi}{2}\,\frac{\bigl(d\sqrt{2/\pi}\bigr)^3}{\bigl(1-2d^2/\pi\bigr)^{3/2}},
\qquad d=\frac{\kappa}{\sqrt{1+\kappa^2}} .
\label{eq:gammakappa}
\end{equation}
To leading order in $\gamma$ the true per-report posterior of $\theta$ is the Gram--Charlier
expansion
\begin{equation}
p(\theta\mid\mu) \;=\; \tfrac{1}{s}\,\varphi\!\Bigl(\tfrac{\theta-\mu}{s}\Bigr)
\Bigl[\,1+\tfrac{\gamma}{6}\,\mathrm{He}_3\!\bigl(\tfrac{\theta-\mu}{s}\bigr)\Bigr]
+O(\gamma^{4/3}),
\qquad \mathrm{He}_3(x)=x^3-3x,
\label{eq:gramcharlier}
\end{equation}
while the \emph{reported} posterior is the Gaussian $\mathcal{N}(\mu,(\tau s)^2)$; the scale $\tau$ is
the single knob on the reported bar, and the skew $\gamma$ is the misspecification. The utility is
piecewise-linear with thresholds $t_1<t_2$ and slopes $k_{\mathrm{over}}$ (over-treatment) and
$k_{\mathrm{under}}$ (under-treatment); the cost asymmetry is $\lambda := k_{\mathrm{under}}/k_{\mathrm{over}}\ge 1$.
In the gap regime $t_1$ is pushed far below $t_2$, so the only active threshold is the treat/escalate
boundary $t_2$ (verified: $\mathrm{P}(\textsc{spare})=0$ exactly).

Two scales calibrate the same bar against different objectives, on independent code paths:
\begin{itemize}
  \item $\tau_{\mathrm{stat}}$ --- \emph{statistical} calibration: the scale at which the reported
  central-$L$ interval covers the truth at the nominal level $L$, i.e.\ the root of
  $C(\tau)=\mathrm{P}(|u|\le z_L\,\tau)=L$ with $z_L=\Phi^{-1}(\tfrac12+\tfrac{L}{2})$ ($L=0.90$).
  % src: minos-core/minos/calibration.py L_ref=0.90 (model constant); RESULTS_B.md.
  It never reads the utility.
  \item $\tau^\star$ --- \emph{decision} calibration: $\tau^\star=\arg\max_\tau \mathbb{E}\,
  U\!\bigl(a^\star(\mathcal{N}(\mu,(\tau s)^2)),\theta\bigr)$, the scale whose induced decision rule
  maximises realised expected utility. It never reads coverage.
\end{itemize}
The \emph{decision--calibration gap} is $G := \tau^\star - \tau_{\mathrm{stat}}$.

\paragraph{The symmetric/degenerate baseline.}
When the posterior is symmetric ($\kappa=0$, so $u\sim\mathcal{N}(0,1)$) or the cost is symmetric
($\lambda=1$), the two calibrations coincide at the same point: both scales return $\approx 1$ and the
gap collapses. At the well-specified cell the simulation gives
$\tau_{\mathrm{stat}}=0.9998$, $\tau^\star=0.9959$, $G=-0.0039$
% src: minos-core/RESULTS_B.md (v2, seed 20240517, RUN_N=4e6), kappa=0 row -- near-degenerate, opposite-tiny sign.
--- a near-degenerate gap, essentially zero. The interesting regime is misspecified and asymmetric.
There the two calibrations fall on \emph{opposite sides of $1$}: statistical calibration
\emph{shrinks} the bar below unity to restore central coverage of a now-skewed truth, while decision
calibration \emph{widens} it above unity so escalation fires soon enough against the heavy
under-treatment tail. At the default cell ($\kappa=3,\lambda=3,\rho=0.5$, so $\gamma=0.6670$
% src: theory/confirm.py, gamma_of_kappa(3).
):
\begin{equation}
\tau_{\mathrm{stat}} \approx 0.964 \;<\; 1 \;<\; \tau^\star \approx 1.05,
\qquad G \approx +0.085 ,
\label{eq:oppositesides}
\end{equation}
% src: theory/confirm.py (A), 5 seeds @ 2e6: tau_stat=0.9640+/-0.0006, tau*=1.0516+/-0.0043, G=0.0876+/-0.0041;
% v2 RESULTS_B: 0.9635 / 1.0431 / 0.0796. Quoted to 2 d.p. as the regime, made exact in Sec.~\ref{sec:thm1}.
a deliberately under-confident bar beats the statistically calibrated one for the decision. The rest
of this section explains \emph{why} this gap exists, \emph{how} it scales, and \emph{what} it implies
for monitoring a correction built on it.

\subsection{Theorem 1: the gap is a scaling law}
\label{sec:thm1}

\begin{theorem}[Decision--calibration gap scaling law]
\label{thm:gap}
In the locked model, to leading order in the posterior skewness $\gamma$,
\begin{align}
\tau_{\mathrm{stat}}(\gamma) &= 1 - a\,\gamma + O(\gamma^2), & a&=0,\\
\tau^\star(\gamma) &= 1 + b\,\Lambda(\lambda)\,\gamma + O(\gamma^{4/3}), & b&=\tfrac16,\quad \Lambda(\lambda)=-z^\star(\lambda)=|z^\star(\lambda)|,\\
G(\gamma) &= \tau^\star-\tau_{\mathrm{stat}} = \tfrac16\,|z^\star(\lambda)|\,\gamma + O(\gamma^{4/3}),
\label{eq:gaplaw}
\end{align}
where $z^\star(\lambda)<0$ is the decision-boundary offset (in reported standard deviations) solving
$(\lambda-1)\,\psi(z)+z=0$ with $\psi(z)=z\,\Phi(z)+\varphi(z)$ the standardised expected
positive part. The gap is positive (the asymmetric cost \emph{widens} the bar,
$\tau^\star>1=\tau_{\mathrm{stat}}$), vanishes as $\gamma\to0$, and is monotone increasing in
$\gamma$ and in the cost asymmetry $\lambda$.
\end{theorem}

\paragraph{Proof sketch (derived symbolically in \texttt{theory/gap\_scaling.py}).}
\emph{(i) $\tau_{\mathrm{stat}}$ is first-order skew-insensitive, $a=0$.} Coverage is a
\emph{symmetric} probability $C(\tau)=\mathrm{P}(|u|\le z_L\tau)$. Its $O(\gamma)$ Gram--Charlier
correction is $\tfrac16\int_{-c}^{c}\varphi(u)\,\mathrm{He}_3(u)\,du$; since $\mathrm{He}_3$ is
\emph{odd} and the interval is symmetric about the mean, the integral is identically zero. Hence the
entire first-order gap is decision-side.
\emph{(ii) $\tau^\star$ slope.} For fixed $\tau$ the reported rule is a pure threshold on $\mu$
(\textsc{escalate} iff $\mu>t_2+\tau s\,z^\star$), so choosing $\tau$ is choosing the threshold. The
realised-utility-optimal threshold solves the first-order condition
$h(z;\gamma)=(\lambda-1)\,g(z;\gamma)+z=0$ with
$g(z;\gamma)=\mathbb{E}_u[(z+u)_+]=\psi(z)+\tfrac{\gamma}{6}\,gc(z)+O(\gamma^2)$, and sympy evaluates
the skew correction in closed form,
\begin{equation}
gc(z)=\int_{-z}^{\infty}(z+u)\,\mathrm{He}_3(u)\,\varphi(u)\,du = -z\,\varphi(z).
\label{eq:gc}
\end{equation}
% src: theory/gap_scaling.py, sympy: gc(z) = -z*phi(z) (asserted equal).
The implicit-function theorem then gives
$\tfrac{d\tau^\star}{d\gamma}\big|_0=\tfrac{(\lambda-1)\varphi(z^\star)}{6[(\lambda-1)\Phi(z^\star)+1]}=b\,\Lambda(\lambda)$
with $b=\tfrac16$; substituting the boundary condition $(\lambda-1)\psi(z^\star)=-z^\star$ collapses
$\Lambda$ to the pure threshold quantity $\Lambda(\lambda)=-z^\star(\lambda)$ (sympy-verified, residual
$0$). Thus $G(\gamma)=\tfrac16|z^\star(\lambda)|\gamma$: the gap's slope in skewness is the magnitude
of the decision-boundary offset, over six. The proximity $\rho$ does \emph{not} enter the leading
coefficient (the optimality condition is pointwise, so the population density cancels); it sets only
whether $\tau^\star$ is identified. At $\lambda=1$, $z^\star=0\Rightarrow\Lambda=0\Rightarrow
\tau^\star=1$: a symmetric cost has no gap. \hfill$\square$

The decision boundary $z^\star(\lambda)$ and the resulting slope are, for $\lambda=1,2,3,4$,
\begin{equation}
z^\star = 0,\,-0.27603,\,-0.43633,\,-0.54916,
\qquad
\tfrac{d\tau^\star}{d\gamma}=\tfrac{|z^\star|}{6}=0,\,0.04600,\,0.07272,\,0.09153 ,
\label{eq:zstargrid}
\end{equation}
% src: theory/gap_scaling.py, GATE 1 grid (mpmath root + sympy slope), this session.
strictly increasing in $\lambda$. At the default asymmetry $\lambda=3$ the leading-order slope is
$|z^\star(3)|/6=0.0727$,
% src: theory/gap_scaling.py / theory/confirm.py (B), slope_th.
giving a leading-order $\tau^\star=1+\tfrac16|z^\star(3)|\gamma=1.0485$ at $\gamma=0.6670$.
% src: theory/gap_scaling.py default-cell print / theory/confirm.py (A), tstar_th = 1.0485.

\paragraph{The leading-order caveat (the operating skew is moderate, not small).}
The default cell sits at $\gamma=0.6670$ --- moderate, \emph{not} small --- so the leading-order law
captures the decision-side scale, sign, and slope, but not the full gap magnitude at that cell. Two
pieces must be combined. The decision side is the leading-order $\tau^\star=1.0485$, which the
simulation confirms to $0.3\%$ ($1.0516\pm0.0043$ over 5 seeds at $2\times10^6$);
% src: theory/confirm.py (A): tau*_theory 1.0485 vs sim 1.0516+/-0.0043, |diff|=0.0031, tol 0.015 PASS.
the calibration side has $a=0$ to first order, but the simulation shows a second-order
\emph{symmetric-coverage shrink} $\tau_{\mathrm{stat}}-1=-0.0360$, reproduced to $0.0002$ by the exact
skew-normal coverage model ($\tau_{\mathrm{stat}}=0.9638$ vs sim $0.9640\pm0.0006$).
% src: theory/confirm.py (A): tau_stat exact 0.9638 vs sim 0.9640+/-0.0006, |diff|=0.0002, tol 0.010 PASS.
Consequently, at $\gamma=0.6670$ the \emph{pure} leading-order decision-side term
$G_{\mathrm{lead}}=\tfrac16|z^\star(3)|\gamma=0.0485$ is only
\begin{equation}
\frac{G_{\mathrm{lead}}}{G_{\mathrm{full}}}=\frac{0.0485}{0.0847}\approx 57\%
\label{eq:57pct}
\end{equation}
% src: theory/confirm.py (A): G_leading 0.0485, G_full (tau*_lead - tau_stat_exact) 0.0847; 0.0485/0.0847=0.573.
of the gap; the remaining $\approx 43\%$ is the higher-order coverage shrink. Combining the two
analytic pieces reconstructs the full gap, $G_{\mathrm{full}}=\tau^\star_{\mathrm{lead}}-
\tau_{\mathrm{stat}}^{\,\mathrm{exact}}=0.0847$, against a simulated $0.0876\pm0.0041$ (and $0.0796$
in v2),
% src: theory/confirm.py (A): |G_full - G_sim| = 0.0029, tol 0.015 PASS; v2 RESULTS_B G=0.0796.
the theory sitting between the two simulation estimates. Across the published skew sweep the full
theory tracks v2 to a maximum deviation of $0.0065$ (tolerance $0.012$); the \emph{pure}
leading-order term tracks the full gap to within $0.012$ up to $\gamma\approx 0.137$, beyond which the
named second-order coverage term must be added.
% src: theory/confirm.py (B): max|G_full - G_v2|=0.0065 < 0.012; lead-order valid up to gamma=0.137.
\textbf{Stated regime of validity:} Theorem~\ref{thm:gap} is a leading-order law in $\gamma$; the
decision-side scale, the gap sign, and the slope hold throughout the tested range, while the full
magnitude at moderate $\gamma$ requires the separately validated second-order coverage shrink in
$\tau_{\mathrm{stat}}$.

The gap is therefore structural, not a numerical accident of one cell: it must appear whenever the
posterior is skewed ($\gamma>0$) and the cost is asymmetric ($\lambda>1$), and its coefficient is the
interpretable threshold offset $z^\star(\lambda)$.

\subsection{Theorem 2: the detectability dichotomy}
\label{sec:thm2}

The gap motivates a \emph{loss-calibration correction} --- fitting the bar's scale to the decision,
i.e.\ reporting $\hat\tau_{\mathrm{cal}}\approx\tau^\star$ rather than $\tau_{\mathrm{stat}}$ (a cited
baseline; see Sec.~\ref{sec:positioning}). Under deployment shift this correction goes stale,
stranding regret
$R(\delta)=\max_\tau \mathrm{EU}(\tau;\delta)-\mathrm{EU}(\hat\tau_{\mathrm{cal}};\delta)\ge 0$.
The deployment shift splits, at matched gain $\beta s$, into two channels that impose the \emph{same}
truth-vs-report discrepancy (hence the same regret) and differ only in whether the observable reported
point moves:
\[
\text{observable } (\delta_{\mathrm{obs}}):\ \mu=\mu_0-\beta s\,\delta_{\mathrm{obs}},
\qquad
\text{hidden } (\delta_{\mathrm{hid}}):\ \theta=\mu_0+\beta s\,\delta_{\mathrm{hid}}+s\,u .
\]
A \emph{label-free monitor} is any measurable $M=f(O)$ of the observable batch $O=\{\mu_i\}$ (plus
fixed known constants); it may not read the truths $\{\theta_i\}$.

\begin{theorem}[Detectability dichotomy]
\label{thm:detect}
Let $P_O^{\mathrm{cal}},P_O^{\mathrm{dep}}$ be the calibration and deployment laws of the observable
reported point $\mu$, and let $M=f(O)$ be any label-free monitor.

\emph{(ii) Achievable bound (observable component).} The stale-correction regret of the observable
channel obeys
\begin{equation}
R_{\mathrm{obs}} \;\le\; L\cdot\delta\bigl(P_O^{\mathrm{dep}},P_O^{\mathrm{cal}}\bigr),
\label{eq:bound}
\end{equation}
in two consistent forms: \textbf{(W$_1$)} $\delta=W_1$ and $L=L_U=k_{\mathrm{under}}$, the utility's
Lipschitz constant in $\theta$; \textbf{(TV)} $\delta=\mathrm{TV}$ and $L=2\,\mathrm{osc}(g_\tau)$,
twice the oscillation of the per-report expected utility
$g_\tau(\mu)=\mathbb{E}[U(a^\star(\mu;\tau),\theta)\mid\mu]$.

\emph{(i) Hidden-channel impossibility.} Define the \emph{hidden} channel as the component of a shift
that leaves $P_O$ exactly invariant (realised by $\delta_{\mathrm{obs}}=0$). Then every label-free
statistic $M=f(O)$ has identical law under \emph{fresh} ($\delta_{\mathrm{hid}}=0$) and \emph{stale}
($\delta_{\mathrm{hid}}>0$) deployment; hence every detector built from $O$ has true-positive rate
equal to its false-positive rate at every operating point, its ROC is the diagonal, and its detection
AUC is exactly $\tfrac12$.
\end{theorem}

\paragraph{Proof sketch of (ii) (\texttt{theory/detectability.py}).}
An optimality chain gives
\[
R(P_{\mathrm{dep}})\le 2\sup_\tau\bigl|\mathrm{EU}(\tau;P_{\mathrm{dep}})-\mathrm{EU}(\tau;P_{\mathrm{cal}})\bigr|=:2\varepsilon .
\]
Writing $\mathrm{EU}(\tau;P)=\int g_\tau\,dP(\mu)$
gives the TV form $\varepsilon\le\mathrm{osc}(g_\tau)\,\mathrm{TV}(P_O)$; and since $U$ is
$L_U$-Lipschitz in $\theta$ while the observable shift only translates the decision boundary by
$\Delta=\beta s\,\delta_{\mathrm{obs}}$ in report-space, the W$_1$ form follows with
$W_1(P_O^{\mathrm{dep}},P_O^{\mathrm{cal}})=\Delta$ exactly. \hfill$\square$

Read off the code's utility, $L_U=\max(k_{\mathrm{under}},k_{\mathrm{over}})=\max(3,1)=3.0$ and the
gain $\beta s=2.5$, so $W_1=2.5\,\delta_{\mathrm{obs}}$; the TV constant is
$L_{\mathrm{TV}}=2\,\mathrm{osc}(g_\tau)=0.6820$.
% src: theory/detectability.py: L_U=3.0, gain=beta*s=2.5, osc(g)=0.3410, L_TV=0.6820, tau_hat_cal=1.0543.
Both bounds hold across the observable sweep $\delta_{\mathrm{obs}}\in[0,0.24]$; the realised regret is
sub-linear (quadratic near $\delta=0$), so the observed regret-vs-$W_1$ slope is $0.3165$, well below
the envelope $L_U=3$.
% src: theory/detectability.py: bound holds across sweep; observed R-vs-W1 slope = 0.3165 <= 3.

\paragraph{Proof of (i) (data-processing / sufficiency; \texttt{theory/impossibility.md}).}
The hidden shift perturbs only $\theta$, so by definition the observable law is exactly invariant,
$P_O^{\mathrm{stale}}=P_O^{\mathrm{fresh}}=:P_O$. Any $M=f(O)$ inherits the invariance,
$\mathrm{Law}(M\mid\mathrm{stale})=\mathrm{Law}(M\mid\mathrm{fresh})$, equivalently
$I(M;\delta_{\mathrm{hid}})\le I(O;\delta_{\mathrm{hid}})=0$; by Neyman--Pearson the likelihood ratio
is $\equiv 1$, so no test has power and $\mathrm{AUC}=\tfrac12$. The statement is finite-sample, not
merely asymptotic: since $M(\mathrm{fresh},\mathrm{seed})=M(\mathrm{stale},\mathrm{seed})$ pathwise,
the fresh and stale score sets coincide, so the empirical AUC is pinned to $\tfrac12$ for every sample
size (a tie at every cross-pair). \hfill$\square$

\textbf{Machine-verified} (\texttt{theory/impossibility\_check.py}, GATE 2(i)): across $16$ seeds and
the full sweep $\delta_{\mathrm{hid}}\in[0,0.24]$ the observable batch and the monitor are bit-identical,
$\max|\Delta\mu|=\max|\Delta M|=0$; the fresh-vs-stale detection AUC is exactly $0.500000$ with a
degenerate bootstrap CI $[0.500000,0.500000]$, and the regret-detection AUC is $0.5000$ (v3: $0.500$).
% src: theory/impossibility_check.py, seeds 8001-8016: max|dmu|=max|dM|=0; fresh-vs-stale AUC=0.500000, CI [0.5,0.5]; regret AUC=0.5000.
Inline, \texttt{detectability.py} (seed 4242) shows the monitor identical at $\delta_{\mathrm{hid}}=0$
versus $0.20$ ($M=0.007248$ in both cases).
% src: theory/detectability.py, seed 4242, 200k voxels: M(delta_hid=0)=M(delta_hid=0.20)=0.007248, points identical.
The result is \emph{conditional only on the definition} of ``hidden'' as the observable-invariant
component, which the Minos model realises by construction; the genuinely observable-moving part of any
real shift falls under part~(ii) (Proposition~\ref{prop:leak}).

\paragraph{Together: a dichotomy.}
The two parts of Theorem~\ref{thm:detect} say a label-free monitor can be made arbitrarily good on the
observable component (regret bounded by a utility-Lipschitz constant times an observable distance) and
is provably powerless on the hidden component, which induces the \emph{same} regret yet leaves every
observable summary untouched. The monitor is not ``imperfect''; it is optimal up to a floor that no
label-free statistic can cross.

\paragraph{Partial leak.}
Real shifts need not factor cleanly into one channel or the other. The following decomposition, a
direct corollary of the dichotomy, states what survives.

\begin{proposition}[Partial leak]
\label{prop:leak}
Decompose a deployment shift into a component that moves the observable law $P_O$ and a residual
component that leaves $P_O$ exactly invariant. Then \emph{(a)} the regret attributable to the
$P_O$-moving (``leaked'') component is bounded-detectable via Theorem~\ref{thm:detect}(ii),
$R_{\mathrm{leak}}\le L_U\,W_1(P_O^{\mathrm{dep}},P_O^{\mathrm{cal}})$; and \emph{(b)} the regret of the
$P_O$-invariant (``genuinely hidden'') component is undetectable by any label-free monitor
($\mathrm{AUC}=\tfrac12$, Theorem~\ref{thm:detect}(i)). A physically ``hidden'' change that correlates
with acquisition features leaks into $P_O$ and is detectable to the extent of that leakage; only the
residual that leaves every observable summary untouched is invisible.
\end{proposition}

\noindent Proposition~\ref{prop:leak}(a) follows from Theorem~\ref{thm:detect}(ii) applied to whatever
moves $P_O$; part~(b) rests on the same definitional exact-invariance premise as
Theorem~\ref{thm:detect}(i) (the \emph{definition} of ``hidden'' in this model).

\subsection{The validity monitor: Theorem 2's applied face}
\label{sec:monitor}

The dichotomy has a direct operational reading. The label-free monitor
\begin{equation}
M \;=\; \sum_b \omega(z_b)\,\bigl|p_{\mathrm{dep}}(z_b)-p_{\mathrm{cal}}(z_b)\bigr|\,dz,
\qquad z=\tfrac{\mu-t_2}{s},\quad \omega(z)=k_{\mathrm{under}}\,\varphi(z),
\label{eq:monitor}
\end{equation}
is a regret-targeted (not density-ratio) divergence of the reported points: $\omega$ weights the
discrepancy by the under-treatment stakes and localises it to within $\sim$one reported standard
deviation of the active threshold $t_2$. By construction $M$ reads only $\{\mu\}$, so it is exactly the
object Theorem~\ref{thm:detect} bounds.

\paragraph{Observable fires; hidden is blind.}
On the observable channel $M$ tracks the stranded regret and separates stale from fresh batches:
$\mathrm{corr}(M,R)=+0.964$ and detection AUC for $\{R>\text{tol}\}$ is $1.000$.
% src: theory/confirm.py (C): corr(M,R)=+0.964, AUC_obs=1.000; v3 RESULTS_C: +0.965, 1.000.
On the hidden channel --- which by construction does not move $\{\mu\}$ --- $M$ sits at the theoretical
floor, detection AUC $=0.500$,
% src: theory/confirm.py (C): AUC_hid=0.500; v3 RESULTS_C: 0.500.
exactly Theorem~\ref{thm:detect}(i). The fitted correction $\hat\tau_{\mathrm{cal}}=1.0543$ (this session;
$1.0480$ in v3, $\tau^\star=1.0485$ in theory) reproduces the decision scale.
% src: theory/confirm.py (C): tau_hat_cal=1.0543; v3 RESULTS_C: 1.0480; theory tau*=1.0485 (Thm 1).

\paragraph{The break-even shift.}
In operation the monitor declares a batch stale when $M$ exceeds a null-calibrated threshold $m^\star$
(false-alarm controlled at a chosen level). The decision-relevant break-even is the shift at which the
stranded regret crosses the tolerance: on the observable channel $R_{\mathrm{obs}}$ crosses the
decision tolerance ($\text{tol}=0.02$) between $\delta_{\mathrm{obs}}=0.06$ ($R=0.0145$) and $0.09$
($R=0.0317$),
% src: theory/confirm.py (C) regret ladder: R_obs at 0.06=0.0145, 0.09=0.0317; TOL=0.02.
and Theorem~\ref{thm:detect}'s envelope $R_{\mathrm{obs}}\le L_U W_1=7.5\,\delta_{\mathrm{obs}}$
guarantees the regret a detectable observable signature wherever it matters. On the hidden channel no
such break-even exists for \emph{any} label-free $M$: the regret ladder is comparable in magnitude
(matched by construction) yet $M$ cannot see it. The monitor is thus the applied face of the
dichotomy --- a near-optimal observable detector paired with a provable blind spot.

\subsection{Positioning: a gap and a monitor, not a new calibration method}
\label{sec:positioning}

Minos's contributions are \emph{(i)} the gap scaling law of Theorem~\ref{thm:gap} --- a quantification
of \emph{how far} the decision-optimal scale departs from the coverage-optimal one --- and \emph{(ii)}
the detectability dichotomy of Theorem~\ref{thm:detect} with the validity monitor as
its applied face. It does \emph{not} propose a new calibration method.

The act of tuning a posterior so the \emph{decision} is optimal is the established baseline of
loss-calibrated Bayesian inference \citep{lacostejulien2011} and its post-hoc form
\citep{vadera2021}, and of decision calibration \citep{zhao2021}; Minos uses exactly this as the
$\hat\tau_{\mathrm{cal}}$ correction and asks two questions these lines do not, since each targets a
single object: how far apart the two calibrations are (the gap), and whether a correction built on one
of them can be monitored for staleness (the dichotomy). The decision-curve / net-benefit literature
\citep{vickers2006} prices a single threshold's net benefit but not the statistical-versus-decision
gap; the calibration-under-shift literature \citep{ovadia2019} studies calibration under distribution
shift empirically but not the label-free detectability floor. To our knowledge Theorem~\ref{thm:detect}(i)
is the first statement that the hidden component of correction-staleness is information-theoretically
invisible to any label-free monitor.

\subsection{The observable/hidden throughline}
\label{sec:throughline}

The two theorems share one structure: a quantity that is decision-relevant yet invisible to the
natural observable. In Theorem~\ref{thm:gap} the skew $\gamma$ moves the decision-optimal scale while
leaving symmetric coverage first-order unchanged ($a=0$); in Theorem~\ref{thm:detect}(i) the hidden shift
moves the truth while leaving every observable summary exactly unchanged. Both motivate reading a
\emph{label}: only a statistic that touches $\theta$ can see what the observables cannot, which is the
role of periodic labeled repeatability spot-checks (Echo, \citealp{echo}). Theorem~\ref{thm:detect}(i) is
the principled justification for pairing a cheap label-free monitor with a sparse labeled check, rather
than trusting either alone.

\paragraph{Impossibility, not identifiability.}
This blind spot must not be confused with an \emph{identifiability} limit. Gauge \citep{gauge}
describes an identifiability wall: a parameter (the IVIM pseudo-diffusion coefficient) that the data
and forward model cannot pin down, so \emph{no estimator} recovers it from the available
measurements. Theorem~\ref{thm:detect}(i) is a different kind of wall: it is an
\emph{information-theoretic impossibility} of \emph{detection} --- the hidden shift contributes zero
mutual information to the observables ($I(M;\delta_{\mathrm{hid}})=0$), so \emph{no label-free test}
can separate stale from fresh, regardless of estimator quality. Identifiability concerns estimating a
parameter from data; this impossibility concerns detecting a distributional change from observables.
Both are ``you cannot get there from here'' statements, but they bound different operations and rest on
different premises --- and only the labeled check (Echo) resolves the impossibility, just as only
additional informative data could resolve an identifiability wall.

% --- References for standalone compile (numbered natbib). In the full manuscript these merge into the
% --- manuscript .bib; keys: lacostejulien2011, vadera2021, zhao2021, vickers2006, ovadia2019, gauge, echo.
